% Source printed pp. 423-425 / PDF pp. 431-433.
\section{Hilbert--Mumford 判据}\label{src:8.4}
\begin{epigraph}
这里先作一点警告, 也致一点歉意. 本文有几千个公式允许一种或多种``类似符号的歧义'': 即另一种对称却不等价的表述. 这些歧义出现在定义和定理中. 我已付出超乎常人的努力\footnote{我没有付出这种超乎常人的努力, 也不保证所有符号都正确.}以求前后一致, 甚至以求陈述正确; 但仍无法保证结果.
\epigraphauthor{David Mumford \cite[第 288 页]{Mum66b}}
\end{epigraph}
\subsection{Hilbert--Mumford 指数}\label{src:8.4.1-subsection}
Hilbert--Mumford 判据 (\bookref{8.4.4}) 常可用于有效计算稳定和半稳定轨迹. 为给出其表述, 设 $U\subseteq\mathbb P(V)$ 是 $G$-等变闭子概形, 其中 $V$ 是有限维 $G$-表示; 这对应于 $U$ 上甚丰沛 $G$-线丛 $L$ (有时称为 $G$-线性化) 的选择. 选择 $V$ 的一组基 $v_1,\ldots,v_n$ 后, 可将点 $p\in\mathbb A(V)(\Bbbk)=V^\vee$ 写成坐标 $(p_1,\ldots,p_n)\in\Bbbk^n$, $p_i=p(v_i)$. 给定单参数子群 $\lambda:\mathbb G_m\to G$, 可选取 $V$ 的基 $v_1,\ldots,v_n$, 使 $\lambda(t)\cdot(p_1,\ldots,p_n)=(t^{d_1}p_1,\ldots,t^{d_n}p_n)$.
\begin{definition}[Hilbert--Mumford 指数]\label{src:8.4.1}
对有提升 $\widetilde u\in\mathbb A(V)$ 的 $\Bbbk$-点 $u\in U$, $u$ 关于 $\lambda$ 的\newterm{Hilbert--Mumford 指数}{Hilbert--Mumford index}为
\begin{equation}\label{src:8.4.2}
\mu(u,\lambda):=\max_{i,\widetilde u_i\ne0}(-d_i).
\end{equation}
有时写作 $\mu_L(u,\lambda)$, 以强调其对丰沛 $G$-线丛 $L$ 的依赖. Hilbert--Mumford 指数还可等价地定义如下: 若 $u_0=\lim_{t\to0}\lambda(t)\cdot u\in\mathbb P(V)$ (因为 $\mathbb P(V)$ 固有, 此极限存在), 则 $\mathbb G_m$ 固定 $u_0$, $\mu_L(u,\lambda)$ 是一维向量空间 $L\otimes\kappa(u_0)$ 上诱导 $\mathbb G_m$-作用的权的相反数.
\end{definition}
\begin{remark}\label{src:8.4.3}
由 Hilbert--Mumford 指数的定义可知:
\begin{enumerate}[(a)]
\item $\lim_{t\to0}\lambda(t)\cdot\widetilde u$ 存在当且仅当 $\mu(u,\lambda)\leq0$;
\item $\lim_{t\to0}t\cdot\widetilde u=0$ 当且仅当 $\mu(u,\lambda)<0$;
\item $\mu(gx,g\lambda g^{-1})=\mu(x,\lambda)$.
\end{enumerate}
\end{remark}
\subsection{GIT 第二基本定理}\label{src:8.4.2-subsection}
\begin{theorem}[Hilbert--Mumford 判据]\label{src:8.4.4}
设代数闭域 $\Bbbk$ 上的线性约化群 $G$ 作用于 $G$-等变闭子概形 $U\subseteq\mathbb P(V)$, 其中 $V$ 是有限维 $G$-表示. 设 $u\in\mathbb P(V)$ 是有提升 $\widetilde u\in\mathbb A(V)$ 的 $\Bbbk$-点. 则
\begin{align*}
u\in U^{\mathrm{ss}}&\quad\Longleftrightarrow\quad0\notin\overline{G\widetilde u}\\
&\quad\Longleftrightarrow\quad\lim_{t\to0}\lambda(t)\cdot\widetilde u\ne0\quad\text{对所有 }\lambda:\mathbb G_m\to G\\
&\quad\Longleftrightarrow\quad\mu(u,\lambda)\geq0\quad\text{对所有 }\lambda:\mathbb G_m\to G.
\end{align*}
若 $G$ 在 $U$ 上的作用一般具有有限稳定子, 则
\begin{align*}
u\in U^{\mathrm s}&\quad\Longleftrightarrow\quad G\widetilde u\subseteq\mathbb A(V)\text{ 闭}\\
&\quad\Longleftrightarrow\quad\mu(u,\lambda)>0\quad\text{对所有非平凡 }\lambda:\mathbb G_m\to G.
\end{align*}
\end{theorem}
\begin{remark}\label{src:8.4.5}
如今称为 Hilbert--Mumford 判据的判据最初由 Hilbert 在 \cite[\S15--16]{Hil93} 中提出, 后由 Mumford 在 \cite[第 53 页]{GIT} 中加以改造. 在 $G$ 约化的更一般情形下, 它也成立.
\end{remark}
\begin{proof}
对半稳定性, 第一个 $(\Rightarrow)$ 蕴含显然: 若 $0\in\overline{G\widetilde u}$, 则对每个非常数不变函数有 $f(\widetilde u)=f(0)=0$, 故 $u\notin U^{\mathrm{ss}}$. 反之, 若 $0\notin\overline{G\widetilde u}$, 则 $0$ 与 $\overline{G\widetilde u}$ 是 $\mathbb A(V)$ 中不交的闭 $G$-不变子概形. 因而它们在 $\mathbb A(V)\sslash G=\Spec(\operatorname{Sym}^*V^\vee)^G$ 中的像不交 (推论~\bookref{7.5.13}(2)). 于是可找到不变函数 $f\in(\operatorname{Sym}^*V^\vee)^G$, 使 $f(0)=0$, $f(\widetilde u)\ne0$, 且可令 $f$ 是正次数齐次元, 即 $f\in\operatorname{Sym}^dV^\vee=\Gamma(\mathbb P(V),\mathcal O(d))$, $d>0$. 第二个等价关系中的 $(\Rightarrow)$ 也显然: 若存在 $\lambda$ 使 $\lim_{t\to0}\lambda(t)\cdot\widetilde u=0$, 则 $0\in\overline{G\widetilde u}$. 反之, 若 $0\in\overline{G\widetilde u}$, 则不稳定化定理 (\bookref{8.3.6}) 给出单参数子群 $\lambda$, 使 $u$ 在 $\lambda$ 下的极限为 $0$. 第三个等价关系由 Hilbert--Mumford 指数的定义得到 (见注记~\bookref{8.4.3}).

对稳定性, 可设 $u\in U^{\mathrm{ss}}$; 否则 $0$ 属于 $G\widetilde u$ 的闭包, 因而 $G\widetilde u$ 不闭. 按定义, 有在 $u$ 处不为零的正次数不变截面 $f\in\Gamma(U,\mathcal O(d))^G$. 必要时增大 $d$ 后, 可使 $f$ 延拓为不变截面 $f\in\Gamma(\mathbb P(V),\mathcal O(d))^G$: 这由正合列 $0\to I_U\to\mathcal O_{\mathbb P(V)}\to\mathcal O_U\to0$, $N\gg0$ 时 $\mathrm H^1(\mathbb P(V),I_U(N))$ 消没以及取不变量的正合性 (即 $G$ 的线性约化性) 得到. 于是可将 $f$ 视为 $\mathbb A(V)$ 上次数 $d$ 的齐次多项式. 令 $\alpha=f(\widetilde u)$, 有交换图
\[
\begin{tikzcd}[ampersand replacement=\&]
G\arrow[r,"\Psi_{\widetilde u}"]\arrow[dr,"\Psi_u"']\&{V(f-\alpha)}\arrow[r,hook]\arrow[d]\&\mathbb A(V)\\
{}\&\mathbb P(V)_f\&{}
\end{tikzcd}
\]
其中 $\Psi_u(g)=g\cdot u$, $\Psi_{\widetilde u}(g)=g\cdot\widetilde u$. 按假设, 有 $\dim G_u=\dim G_{\widetilde u}=0$, 故两个稳定子均有限, 因而固有. 由习题~\bookref{4.8.6}(b), $Gu\subseteq\mathbb P(V)_f$ (分别地, $G\widetilde u\subseteq\mathbb A(V)$) 闭当且仅当 $\Psi_u$ (分别地, $\Psi_{\widetilde u}$) 固有. 另一方面, $V(f-\alpha)\to\mathbb P(V)_f$ 固有, 故 $\Psi_u$ 固有当且仅当 $\Psi_{\widetilde u}$ 固有. 于是 $Gu\subseteq U_f$ 闭当且仅当 $G\widetilde u\subseteq\mathbb A(V)$ 闭, 得到第一个等价关系. 对第二个等价关系, 若 $G\widetilde u$ 不闭, 则存在单参数子群 $\lambda:\mathbb G_m\to G$, 使 $\lim_{t\to0}\lambda(t)\cdot\widetilde u$ 存在且不属于 $G\widetilde u$. 这给出满足 $\mu(u,\lambda)\leq0$ 的非平凡 $\lambda$. 反之, 若 $G\widetilde u$ 闭, 则 $\Psi_{\widetilde u}$ 固有, 故对每个非平凡 $\lambda$, 由 $t\mapsto\lambda(t)\cdot\widetilde u$ 定义的映射 $\mathbb G_m\to\mathbb A(V)$ 也固有; 这意味着 $\lim_{t\to0}\lambda(t)\widetilde u$ 不存在, 因而 $\mu(u,\lambda)>0$.
\end{proof}
\begin{remark}[其他表述]\label{src:8.4.6}
给定 $u\in\mathbb P(V)$ 与单参数子群 $\lambda:\mathbb G_m\to G$, 若 $u$ 关于诱导 $\mathbb G_m$-作用半稳定, 则称 $u$ 为 $\lambda$-半稳定. 用此术语, Hilbert--Mumford 判据断言 $u$ 为 $G$-半稳定当且仅当它对所有 $\lambda$ 均为 $\lambda$-半稳定.

关于环面 $T=\mathbb G_m^k$ 的半稳定性有很好的组合表述. 设 $V=\bigoplus_\chi V_\chi$ 是按特征标 $\chi\in\mathbb X^*(T)\cong\mathbb Z^k$ 索引的权空间分解. 对 $u$ 的提升 $\widetilde u\in V$, 以 $\widetilde u_\chi\in V_\chi$ 表示其投影. $u$ 的\newterm{状态多胞形}{state polytope}为
\[
\operatorname{StatePoly}_T(u):=\operatorname{conv}\{\chi\in\mathbb X^*(T)\mid\widetilde u_\chi\ne0\}
\subseteq\mathbb X^*(T)_{\mathbb Q}\cong\mathbb Q^k.
\]
则 $u$ 为 $T$-半稳定 (分别地, $T$-稳定), 当且仅当平凡特征标 $0\in\operatorname{StatePoly}_T(u)$ (分别地, $0$ 属于其内部).
\end{remark}
\subsection{Hilbert--Mumford 判据的变体}\label{src:8.4.3-subsection}
Hilbert--Mumford 判据 (\bookref{8.4.4}) 可化为判断点 $u\in[U/G]$ 是否半稳定的叠论判据, 即判断 $u$ 是否属于开子叠 $[U^{\mathrm{ss}}/G]$. $G$-等变嵌入 $U\subseteq\mathbb P(V)$ 对应于 $[U/G]$ 上线丛 $L$, 其沿 $U\to[U/G]$ 的拉回甚丰沛. 因为稳定和半稳定轨迹 $G$-不变, 它们定义 $[U/G]$ 的开子叠. 一个点 $u\in U(\Bbbk)$ 和一个单参数子群 $\lambda:\mathbb G_m\to G$ 的数据, 在共轭意义下, 由映射 $f_{u,\lambda}:[\mathbb A^1/\mathbb G_m]\to[U/G]$ 分类, 使诱导映射
\[
\begin{tikzcd}[ampersand replacement=\&]
B\mathbb G_m\arrow[r,"0"]\&{[\mathbb A^1/\mathbb G_m]}\arrow[r,"{f_{u,\lambda}}"]\&{[U/G]}\arrow[r]\&BG
\end{tikzcd}
\]
对应于 $\lambda$ (见命题~\bookref{7.9.1}). Hilbert--Mumford 指数为 $\mu(u,\lambda)=-\mathrm{wt}(f_{u,\lambda}^*L)|_{B\mathbb G_m}$.
\begin{corollary}[Hilbert--Mumford 判据]\label{src:8.4.7}
设代数闭域 $\Bbbk$ 上的线性约化群 $G$ 作用于射影 $\Bbbk$-概形 $U$. 设 $L$ 是 $[U/G]$ 上拉回到 $U$ 后甚丰沛的线丛. 则 $u\in[U/G]$ 半稳定当且仅当对所有映射
\[
f:[\mathbb A^1/\mathbb G_m]\longrightarrow[U/G],\qquad f(1)\simeq u,
\]
均有 $\mathrm{wt}((f^*L)|_{B\mathbb G_m})\geq0$.

若 $G$ 在 $U$ 上的作用一般具有有限稳定子, 则 $u$ 稳定当且仅当对所有满足 $f(1)\simeq u$ 且诱导稳定子映射 $\mathbb G_m\to G_{f(0)}$ 非平凡的映射 $f:[\mathbb A^1/\mathbb G_m]\to[U/G]$, 均有 $\mathrm{wt}((f^*L)|_{B\mathbb G_m})>0$.
\end{corollary}
\begin{remark}[\texorpdfstring{$\Theta$}{Θ}-半稳定性]\label{src:8.4.8}
上述表述提示我们将半稳定性自然推广到任意配有线丛 $\mathcal L$ 的代数叠 $\mathcal X$. 令 $\Theta=[\mathbb A^1/\mathbb G_m]$. 若对所有满足 $f(1)\simeq x$ 的映射 $f:\Theta\to\mathcal X$, 均有 $\mathrm{wt}(f^*\mathcal L)|_{B\mathbb G_m}\geq0$, 则称点 $x\in\mathcal X(\Bbbk)$ 为\newterm{$\Theta$-半稳定}{$\Theta$-semistable}.
\end{remark}
\begin{exercise}[仿射 Hilbert--Mumford 判据]\label{src:8.4.9}
设代数闭域 $\Bbbk$ 上的线性约化群 $G$ 作用于有限型仿射概形 $U=\Spec A$. 设 $\chi:G\to\mathbb G_m$ 是特征标, $U^{\mathrm{ss}}$ 与 $U^{\mathrm s}$ 是习题~\bookref{8.2.10} 定义的关于 $\chi$ 的半稳定和稳定轨迹. 对 $u\in U(\Bbbk)$, 证明
\[
\begin{aligned}
u\in U^{\mathrm{ss}}\quad\Longleftrightarrow\quad&\text{对所有使 }\lim_{t\to0}\lambda(t)\cdot u\text{ 存在的}\\
&\text{单参数子群 }\lambda:\mathbb G_m\to G,\ \langle\chi,\lambda\rangle\geq0,
\end{aligned}
\]
其中 $\langle-,-\rangle$ 为特征标与单参数子群之间的自然配对. 若 $G$ 在 $U$ 上的作用一般具有有限稳定子, 证明 $u\in U^{\mathrm s}$ 当且仅当上述条件以严格不等式 $\langle\chi,\lambda\rangle>0$ 成立.

\textit{提示:} 考虑 $\chi$ 诱导的 $G$ 在 $U\times\mathbb A^1$ 上的作用 $g\cdot(u,z)=(g\cdot u,\chi(g)^{-1}\cdot z)$, 证明 $u\notin U^{\mathrm{ss}}$ 当且仅当 $\overline{G\cdot(u,1)}\cap(U\times\{0\})\ne\emptyset$. 用不稳定化定理 (\bookref{8.3.6}) 证明这等价于存在单参数子群 $\lambda$, 使
\[
\lim_{t\to0}\lambda(t)\cdot(u,1)=\lim_{t\to0}(\lambda(t)\cdot u,t^{-\langle\chi,\lambda\rangle})\in U\times\{0\}.
\]
\end{exercise}
